% 人工对照底本逐节点转写；独立TikZ重绘。

\begin{tikzpicture}[x=1cm,y=-1cm,every node/.style={anchor=west,align=left,inner sep=2pt,font=\fontsize{9.5}{12.3}\selectfont},edge/.style={draw=black!65,line width=.45pt}]
\node[text width=0.788cm,draw=black!65] (n1) at (0.000,1.056) {行为};
\node[text width=0.788cm,draw=black!65] (n2) at (2.120,0.275) {作为};
\node[text width=2.791cm,draw=black!65] (n3) at (4.840,0.275) {“持有”属于作为};
\node[text width=1.122cm,draw=black!65] (n4) at (2.120,1.446) {不作为};
\node[text width=2.123cm,draw=black!65] (n5) at (4.840,1.056) {真正不作为犯};
\node[text width=2.457cm,draw=black!65] (n6) at (4.840,1.837) {不真正不作为犯};
\draw[edge] (n2.east)--(4.690,0.275)--(4.690,0.275)--(n3.west);
\draw[edge] (n1.east)--(1.970,1.056)--(1.970,0.275)--(n2.west);
\draw[edge] (n4.east)--(4.690,1.446)--(4.690,1.056)--(n5.west);
\draw[edge] (n4.east)--(4.690,1.446)--(4.690,1.837)--(n6.west);
\draw[edge] (n1.east)--(1.970,1.056)--(1.970,1.446)--(n4.west);
\end{tikzpicture}

